% English counterpart of draft/background_and_motivation.md.
% Compile from paper/; see motivation-preview.tex for EuroSys 2027 settings.
\section{Background and Motivation}
\label{sec:background-motivation}

Multi-agent applications complete a user task through interdependent LLM
requests. After a worker failure, rebuilding lost key--value (KV) state can
restore an individual request, but the user-visible penalty depends on when
the entire workflow completes. Workflow dependencies and the cost of
rebuilding long contexts motivate allocating limited KV protection resources
to reduce workflow completion time (WCT).

\input{sections/motivation-figures}

\subsection{Workflows, Failures, and KV Recovery}
\label{sec:motivation-background}

\paragraph{Multi-agent workflows.}
We represent a workflow as a directed acyclic graph $G=(V,E)$, whose nodes
are LLM requests and whose edges are data dependencies. An agent branch may
include multiple requests and tool calls; a successor runs only after its
dependencies are satisfied. WCT spans the release of the first request to
the completion of the last. For a given execution schedule, a request's
\emph{slack} is the maximum delay to its completion that does not postpone
workflow completion. Requests on a critical path have zero slack.

\paragraph{Serving and failure model.}
We consider a cluster of GPU workers with block-managed KV that can be
replicated to other GPUs or a host DRAM pool, consistent with KV-centric
architectures such as Mooncake~\cite{qin2025mooncake}. Our failure model is
a fail-stop failure of one GPU or worker: its local KV becomes unavailable,
while the control plane, backup DRAM, and other workers survive. Surviving
workers with the same model already loaded take over affected requests.
Whole-host failures, network partitions, and silent corruption are outside
this failure model.

\paragraph{Recovery paths and resource costs.}
Without a usable copy, the system must rerun prefill to rebuild lost KV.
A DRAM copy avoids reconstructing the saved prefix but must be transferred
to a GPU before use; concurrent recoveries share transfer bandwidth.
A copy already resident in the actual takeover GPU's HBM can be used
directly, at the cost of HBM capacity that could serve foreground requests
and replication traffic before failure. Backup placement therefore trades
capacity and normal-service overhead against recovery speed.

Prior systems reduce failure costs through KV streaming in
D\'{e}j\`{a}Vu~\cite{strati2024dejavu}, erasure coding in
GhostServe~\cite{jayakody2026ghostserve}, and load-aware checkpoint and
recovery coordination in LUMEN~\cite{cao2026lumen}. Multi-agent
applications add an allocation problem: saving the same recovery time for
different requests can have different value to the user. Neither saved KV
bytes nor avoided computation alone captures that difference.

\subsection{Recovery Cost Is Not Workflow Delay}
\label{sec:motivation-workflow}

A join waits for its last required branch. Earlier branches can tolerate
some delay, whereas delaying the branch that determines the join time
postpones downstream work. To characterize branch timing differences, we
analyze the public CC/Weka coding-agent
traces~\cite{semianalysis2026cctraces}. From 393 sessions, we first extract
strictly overlapping subagent groups: the latest branch start precedes
the earliest branch finish. We retain groups whose branches all complete
and are followed by a Main request after the last branch finishes.
Main may also issue requests while branches execute. Because the trace
does not record result dependencies, this criterion infers fan-in from
timing; it does not verify consumption of every branch's output.

The filter retains 286 groups across 121 workflows. For each retained
group, we compute the longest minus the shortest branch duration; each
workflow contributes its maximum such gap.
Figure~\ref{fig:motivation}(\subref{fig:motivation-branch-gap}) shows a median of
138\,s and a 90th percentile of 1,173\,s. These gaps establish substantial
branch-duration heterogeneity. With unequal branch start times, however,
a duration gap is neither join slack nor a direct measure of fault tolerance.

Figure~\ref{fig:motivation}(\subref{fig:motivation-workflow}) shows a
real ten-branch group. Its durations range from 24.960 to 193.273\,s,
with a gap of 168.313\,s, close to the population median of 137.758\,s.
This group contributes its workflow's maximum to
Figure~\ref{fig:motivation}(\subref{fig:motivation-branch-gap}). Nine branches
last less than 117\,s each. A10 finishes last at 221.8\,s after the fork,
followed by a Main request at 226.6\,s. Bars aggregate each branch's
requests and tool calls while preserving observed starts and durations.
Assuming that the successor depends on all ten branches, each branch's
slack is its time margin to A10's completion. We select this example to
illustrate greater fan-out and recovery exposure, not average failure impact.

Figure~\ref{fig:motivation}(\subref{fig:motivation-recovery-exposure}) shows
ten independent counterfactuals at 36.1\,s after the fork, each losing
one branch's history KV. All ten branches have active LLM requests at
that instant. Among such intervals, we choose one with a longer A10
input to illustrate a larger recovery penalty. Trace hash-block counts
provide history lengths $M_i$ of 2,944--123,648 tokens, to which we add an
assumed known 64-token tail. We query the Qwen3-14B configuration of
Figure~\ref{fig:motivation-recompute} at each exact length: $T_i$ estimates
empty-cache reconstruction of $M_i+64$ tokens on A100-SXM, with BF16,
batch one, and TP one. The tail state is a counterfactual assumption,
not a trace observation. Slack $\sigma_i$ comes from trace completion
times. Holding other branches, downstream execution, and queueing
unchanged, the WCT penalty is $\max(0,T_i-\sigma_i)$.
This cost substitution does not constitute physical failure replay on
the original serving model.

A5 has the highest rebuilding cost, 45.7\,s, but its 146.1\,s slack
absorbs the entire delay. A10 costs less to rebuild, at 27.5\,s, yet
exposes that entire delay to WCT because its slack is zero. All other
branches absorb their rebuilding costs. Thus, avoiding the most
recomputation need not yield the largest WCT improvement. KV protection
must account for delay tolerance within a workflow; online estimation and
resource contention require the mechanisms and evaluation developed later.

\subsection{Long Contexts and Recovery Tiers}
\label{sec:motivation-long-context}

\input{sections/motivation-prefill-figure}

Repeated agent interactions accumulate code, tool outputs, and other
agents' intermediate results. Saving history KV avoids repeating that
computation after failure, but its location still affects recovery
latency. Figure~\ref{fig:motivation-recompute} compares three initial
states: the entire history is resident in the takeover GPU's HBM, the
same state is available only in host-local DRAM, or no copy is available.
All paths end after the same known 64-token tail's forward pass, with
history and tail KV available for continued decoding; model weights are
already loaded. HBM reuse therefore still incurs tail computation.

We use AIConfigurator~\cite{nvidia2026aiconfigurator} 0.9.0 HYBRID static
estimates for Qwen3-4B, 8B, and 14B, with A100-SXM, vLLM 0.14.0, BF16,
batch one, and TP one throughout. HYBRID combines available kernel
performance tables with model-based fallbacks; these request-level
latencies are estimates, not physical recovery measurements. HBM reuse
sets the cached prefix to $M$ and computes 64 new tokens; Recompute
starts empty and processes $M+64$ tokens. DRAM reuse first transfers the
entire history, then performs the same tail computation. We calculate
KV bytes from each model's layers, KV heads, and head dimension, and use
the system configuration's assumed 32\,GB/s unidirectional PCIe bandwidth.
This serialized transfer model assumes no transfer/compute overlap.

At 128K history tokens, HBM reuse costs an estimated 35--59\,ms across
the three models, DRAM reuse 639--730\,ms, and Recompute 31.4--49.6\,s.
For Qwen3-14B, a DRAM copy avoids most historical recomputation; full
HBM residency additionally removes about 671\,ms of modeled transfer
waiting. Backups therefore avoid computation, while HBM's incremental
value over DRAM comes from reducing transfer waiting. The magnitude
requires matched physical measurements and concurrent-recovery evaluation.

Full-history HBM residency is a reference, not AegisKV's limited hot-block
peer placement. These shape estimates do not validate long-context
accuracy or model resource contention, and do not establish WCT gains
under a limited backup budget. They motivate allocating protection
between avoiding recomputation and reducing transfer waiting.
Figure~\ref{fig:motivation} further shows why such savings must be
interpreted in the context of branch delay tolerance.

\subsection{Problem and Design Goals}
\label{sec:motivation-goals}

AegisKV asks: \emph{with limited backup capacity and replication and
recovery bandwidth, which KV from which requests should receive what
protection to reduce the workflow-level impact of failures?}
Dependencies differentiate protection benefits, long contexts increase
state-loss costs, and replication competes with foreground serving.

This problem imposes five goals and constraints.
\textbf{(R1)} Optimize recovery for WCT, accounting for dependency-induced
delay tolerance.
\textbf{(R2)} Allocate limited backup capacity across requests and their
KV to reduce exposed recovery waiting.
\textbf{(R3)} Support concurrent recovery and transfer--compute overlap
without introducing request-level serialization for protection priority.
\textbf{(R4)} Bound normal-service replication overhead and subordinate
backup capacity and bandwidth use to foreground needs.
\textbf{(R5)} Preserve model computation and KV recovery correctness using
existing serving resources.
The next section develops the evidence for mechanism choices and the
corresponding placement and recovery design.
